% Full derivation of the condition for cluster-approximated attention.
% Self-contained section: objective, notation, average approximation, the
% cosh condition (original, with corrected proofs), the variance (Bennett)
% refinement, the per-cluster guarantee, and the hybrid per-step certificate.
% Requires: tcolorbox, amsmath.

\section{Approximation}

\paragraph{Goal}
Suppose we split the KVs into clusters with close keys. Under what condition can we
faithfully approximate full attention? We need conditions that can be efficiently
verified, and we need to bound the approximation error.

\begin{tcolorbox}[colback=white, colframe=black, arc=0mm, boxrule=0.5pt]
For a query $q$ and a KV set clustered as $\mathcal C = \{C_1,\dots,C_m\}$: for all
$\epsilon > 0$, using an approximation method $\mathbf{A}$ [efficient], under a condition
$\mathbf{Cond}$ [depending on $q$ and $\epsilon$, efficiently verifiable], the vanilla
output $o$ and the approximated output $\hat o$ satisfy $\|\hat o - o\|_2 < \epsilon$.
\end{tcolorbox}

Different approximation methods $\mathbf{A}$ lead to different conditions $\mathbf{Cond}$.
We use the \emph{average approximation}: each cluster is represented by its mean key and
mean value $(\bar k_C, \bar v_C)$.

\subsection{Notation and setup}

The visible tokens are partitioned into clusters $C_1, \dots, C_m$. For a cluster $C$
let $\bar k_C = \frac{1}{|C|}\sum_{i \in C} k_i$ and
$\bar v_C = \frac{1}{|C|}\sum_{i \in C} v_i$ denote the \emph{arithmetic} means (we drop
the subscript on $\bar k, \bar v$ when the cluster is clear). Define
\[
s_i = \frac{q k_i}{\sqrt{d_k}}, \qquad
s_C = \frac{q \bar k}{\sqrt{d_k}}, \qquad
Z_C = \sum_{i\in C} e^{s_i}, \qquad
\hat Z_C = |C|\, e^{s_C}, \qquad
u_C = \frac{\sum_{i\in C} e^{s_i} v_i}{Z_C},
\]
\[
p_C = \frac{Z_C}{\sum_{C'} Z_{C'}}, \qquad
\hat p_C = \frac{\hat Z_C}{\sum_{C'} \hat Z_{C'}} .
\]
Vanilla attention output and its average approximation:
\[
o = \sum_C p_C\, u_C, \qquad\qquad \hat o = \sum_C \hat p_C\, \bar v_C .
\]

Per-cluster statistics (each is a prompt-side summary plus $O(d_k)$ query-time work):
\begin{align*}
B    &\ \ge\ \max_i \|v_i\|_2 \quad \text{(over all visible tokens)}, \\
B_C  &\ \ge\ \max_{i\in C} \|v_i\|_2, \\
\delta_C &\ \ge\ \max_{i\in C} \Big|\frac{q(k_i - \bar k)}{\sqrt{d_k}}\Big|, \\
\sigma_C^2 &\ \ge\ \frac{1}{|C|}\sum_{i\in C}(s_i - s_C)^2
      \;=\; \frac{q\, \Sigma_C\, q^\top}{d_k},
      \qquad \Sigma_C = \frac{1}{|C|}\sum_{i\in C}(k_i-\bar k)(k_i-\bar k)^\top .
\end{align*}

\emph{Remark (mean-zero property).} Because $\bar k$ is the arithmetic mean of the
cluster keys, $\sum_{i \in C}(s_i - s_C) = 0$ exactly. The whole mass-term analysis
below relies on this; replacing $\bar k$ by a medoid or a normalized centroid breaks
it. It also gives $\sigma_C \le \delta_C$ when both are computed exactly.

\subsection{Error decomposition}

Define the Jensen gap of cluster $C$ and the global inflation factor
\[
b_C = \log \frac{Z_C}{\hat Z_C}
    = \log\Big(\frac{1}{|C|}\sum_{i\in C} e^{s_i - s_C}\Big) \ \ge\ 0,
\qquad
A = \sum_{C'} \hat p_{C'}\, e^{b_{C'}} \ \ge\ 1,
\]
where $b_C \ge 0$ follows from Jensen's inequality and the mean-zero property
(equality iff the scores in $C$ are constant). By construction
$p_C = \hat p_C\, e^{b_C} / A$.

From $\hat o - o = \sum_C (\hat p_C - p_C)\, u_C + \sum_C \hat p_C (\bar v_C - u_C)$
and $\|u_C\|_2 \le B$ ($u_C$ is a convex combination of the $v_i$):
\[
\|\hat o - o\|_2
\ \le\
\underbrace{B \sum_C |p_C - \hat p_C|}_{\text{mass term}}
\; + \;
\underbrace{\sum_C \hat p_C\, \|\bar v_C - u_C\|_2}_{\text{value term}} .
\]
The two terms are the two ways a cluster summary can fail: $\hat Z_C$ mis-estimates the
cluster's softmax mass (spread of scores), and $\bar v_C$ ignores that attention inside
the cluster is softmax-weighted rather than uniform (key--value correlation).

\subsection{Mass term}

\paragraph{Step 1: cancellation across clusters}
Treating $e^{b_C}$ as a random variable with law $\hat p$ (so that its mean is $A$),
and using $\mathbb E|X - \mathbb E X| = 2\,\mathbb E (X - \mathbb E X)_+$,
\begin{align*}
\sum_C |p_C - \hat p_C|
  &= \sum_C \hat p_C \Big| \frac{e^{b_C}}{A} - 1 \Big|
   = \frac{1}{A} \sum_C \hat p_C \big| e^{b_C} - A \big|
   = \frac{2}{A} \sum_C \hat p_C \big( e^{b_C} - A \big)_+ \\
  &\le \frac{2}{A} \sum_C \hat p_C \big( e^{b_C} - 1 \big)
   = \frac{2(A-1)}{A},
\end{align*}
where the inequality uses $A \ge 1$ and $e^{b_C} \ge 1$. Since $A \mapsto 2(A-1)/A$ is
increasing, any upper bound on $e^{b_C}$ yields a bound on the mass term.

\paragraph{Step 2: bounding the Jensen gap}

\begin{tcolorbox}[colback=white, colframe=black, arc=0mm, boxrule=0.5pt]
\textbf{Lemma 1 (variance-constrained bound; Bennett).}
Let $x_i = s_i - s_C$, so $\frac{1}{|C|}\sum_i x_i = 0$,
$\frac{1}{|C|}\sum_i x_i^2 \le \sigma_C^2$, and $x_i \le \delta_C$. Then
\[
e^{b_C} \;=\; \frac{1}{|C|}\sum_{i \in C} e^{x_i}
\;\le\;
G(\sigma_C, \delta_C)
\;:=\;
\frac{\sigma_C^2\, e^{\delta_C} + \delta_C^2\, e^{-\sigma_C^2/\delta_C}}
     {\sigma_C^2 + \delta_C^2} .
\]
\end{tcolorbox}

\emph{Proof.} Let $g(t) = (e^t - 1 - t)/t^2$ (with $g(0) = \tfrac12$); $g$ is
nondecreasing. For any $a$ and any $x \le \delta$,
\[
e^{x} = e^{a}\big[1 + (x-a) + g(x-a)(x-a)^2\big]
\ \le\ e^{a}\big[1 + (x-a) + g(\delta-a)(x-a)^2\big].
\]
Averaging over $i \in C$ and using the zero mean and the second-moment bound,
\[
\frac{1}{|C|}\sum_i e^{x_i}
\ \le\ e^{a}\big[1 - a + g(\delta-a)\,(\sigma^2 + a^2)\big].
\]
Choosing $a = -\sigma^2/\delta$ (so $\delta - a = (\delta^2+\sigma^2)/\delta$ and
$\sigma^2 + a^2 = \sigma^2(\delta^2+\sigma^2)/\delta^2$) and simplifying gives exactly
$G(\sigma,\delta)$. $\hfill\square$

Remarks:
\begin{itemize}
\item[(i)] \emph{The chord bound is the special case with no variance information:}
  $G(\delta, \delta) = \cosh\delta$, recovering
  $e^{b_C} \le \tfrac12(e^{-\delta_C} + e^{\delta_C}) = \cosh\delta_C$, which follows
  directly from the chord inequality
  $e^x \le \frac{\delta - x}{2\delta} e^{-\delta} + \frac{\delta + x}{2\delta} e^{\delta}$
  averaged with mean zero. Also $G(0,\delta) = 1$.
\item[(ii)] Only the one-sided bound $x_i \le \delta_C$ is used, and the bound is tight
  (equality for scores supported on $\{-\sigma^2/\delta,\ \delta\}$).
\item[(iii)] The proof holds verbatim with any over-estimates
  $\hat\sigma_C \ge \sigma_C$, $\hat\delta_C \ge \delta_C$ (choose
  $a = -\hat\sigma_C^2/\hat\delta_C$), so efficiently verifiable upper bounds may be
  plugged in; one may always use $\min\{\cosh\delta_C,\ G(\hat\sigma_C,\hat\delta_C)\}$,
  so the refined condition is never worse than the cosh one.
\end{itemize}

Writing $F_C$ for either $\cosh\delta_C$ or $G(\sigma_C,\delta_C)$ and
\[
S_F = \sum_C \hat p_C\, F_C \ \ge\ A,
\qquad\text{we conclude}\qquad
\sum_C |p_C - \hat p_C| \ \le\ \frac{2\,(S_F - 1)}{S_F}.
\]

\subsection{Value term}

\begin{tcolorbox}[colback=white, colframe=black, arc=0mm, boxrule=0.5pt]
\textbf{Lemma 2 (softmax vs.\ uniform).}
Let $w_i = e^{s_i} / \sum_{j\in C} e^{s_j}$ for $i \in C$ with
$\max_{i\in C}|s_i - s_C| \le \delta_C$. Then
\[
\sum_{i\in C} \Big| \frac{1}{|C|} - w_i \Big| \ \le\ 2 \tanh(\delta_C/2).
\]
\end{tcolorbox}

\emph{Proof.} Let $Y_i = e^{s_i - s_C} \in [a, b]$ with $a = e^{-\delta}$,
$b = e^{\delta}$, and let $T = \frac{1}{|C|}\sum_i Y_i$. The left side equals
$\frac{1}{|C|}\sum_i |Y_i - T| \,/\, T$. For $y \in [a,b]$ the convex function
$y \mapsto (y - T)_+$ lies below its chord, $(y-T)_+ \le (b-T)\frac{y-a}{b-a}$, so
averaging gives $\frac{1}{|C|}\sum_i (Y_i - T)_+ \le \frac{(b-T)(T-a)}{b-a}$, and with
$\frac{1}{|C|}\sum_i |Y_i - T| = \frac{2}{|C|}\sum_i (Y_i - T)_+$,
\[
\sum_{i} \Big| \frac{1}{|C|} - w_i \Big|
\ \le\ \frac{2\,(b-T)(T-a)}{(b-a)\,T} .
\]
The right side is maximized over $T$ at $T = \sqrt{ab} = 1$, where it equals
$\frac{2(b-1)(1-a)}{b-a} = \frac{2(2\cosh\delta - 2)}{2\sinh\delta} = 2\tanh(\delta/2)$.
$\hfill\square$

Since $\sum_i \big(\frac{1}{|C|} - w_i\big) = 0$, for any fixed vector $c$,
\[
\bar v_C - u_C = \sum_{i\in C} \Big(\frac{1}{|C|} - w_i\Big) v_i
             = \sum_{i\in C} \Big(\frac{1}{|C|} - w_i\Big) (v_i - c),
\qquad\text{hence}\qquad
\|\bar v_C - u_C\|_2 \ \le\ 2 B_C \tanh(\delta_C/2)
\]
with $c = 0$ and $B_C = \max_{i\in C}\|v_i\|_2$. (The centered choice
$c = \bar v_C$, $B_C = \max_i \|v_i - \bar v_C\|_2$, is also valid and never larger.)

\subsection{The condition}

Combining the two terms, for $F_C \in \{\cosh\delta_C,\ G(\sigma_C,\delta_C)\}$:

\begin{tcolorbox}[colback=white, colframe=black, arc=0mm, boxrule=0.5pt]
\[
\sum_{C} \hat p_{C}\left(
\frac{2B\,\big(F_C - 1\big)}{S_F}
+ 2 B_C \tanh(\delta_C/2)\right) \ <\ \epsilon
\quad\Longrightarrow\quad
\|\hat o - o\|_2 \ <\ \epsilon,
\qquad S_F = \sum_{C} \hat p_C F_C .
\]
\end{tcolorbox}

With $F_C = \cosh\delta_C$ this is the original condition
(note $\sum_C \hat p_C\, 2B(\cosh\delta_C - 1)/S_F = 2B(S_F-1)/S_F$ since
$\sum_C \hat p_C = 1$). With $F_C = G(\sigma_C, \delta_C)$ it is strictly tighter.

\emph{Magnitude.} Empirically $b_C \approx \sigma_C^2/2$ (rank correlation $0.95$),
with typical $\sigma_C \approx 1$ while the range bound gives
$\delta_C \approx 10\text{--}17$; there $G \approx 2\times 10^2\text{--}8\times 10^4$
versus $\cosh\delta_C \approx 10^4\text{--}10^7$, i.e.\ the per-cluster mass factor
tightens by roughly two orders of magnitude. (Caveat: once $S_F \gg 1$ the
\emph{normalized} mass term saturates at $2B$ for either $F$; the refinement matters
for per-cluster scores, selection quality, and the certificate below once the residual
mass is small.)

\subsection{Per-cluster guarantee}

The summands of the condition are a \emph{diagnostic decomposition}: only their sum
bounds the error. A genuine per-cluster guarantee follows from bracketing both sides of
$|p_C - \hat p_C| = \hat p_C\, |e^{b_C} - A| / A$ with $e^{b_C} \in [1, F_C]$ and
$A \in [1, S_F]$:
\[
\big\| p_C u_C - \hat p_C \bar v_C \big\|_2
\ \le\
\hat p_C \left[\, B\, \max\!\Big( F_C - 1,\ \frac{S_F - 1}{S_F} \Big)
+ 2 B_C \tanh(\delta_C/2) \right].
\]
The second branch of the $\max$ is why the diagnostic share is not a per-cluster bound:
a cluster with $\delta_C = 0$ still carries probability error because \emph{other}
clusters inflate $A$. Summing this guarantee loses the cancellation of Step~1 and is
weaker than the condition by roughly a factor $S_F$; use it per cluster only.

\subsection{Hybrid computation and per-step certificate}

The decode kernel does not approximate every cluster: a selected set $S$ is computed
exactly \emph{inside the same softmax}, i.e.\ for $C \in S$ the mass is the true $Z_C$
and the value is $u_C$ (so $b_C = 0$), while $C \notin S$ keep
$(\hat Z_C, \bar v_C)$. Let $\hat o_{\mathrm{hyb}}$ denote this output and let all
$\hat p_C$ below be the \emph{all-approximate} weights (available before selection).

\begin{tcolorbox}[colback=white, colframe=black, arc=0mm, boxrule=0.5pt]
For any selected set $S$, with
$T = \sum_{C \notin S} \hat p_C\, (F_C - 1)$:
\[
\|\hat o_{\mathrm{hyb}} - o\|_2
\ \le\
\frac{2B\,T}{1+T} \;+\; \sum_{C\notin S} 2\,\hat p_C B_C \tanh(\delta_C/2)
\ \le\
\sum_{C \notin S} \hat p_C \Big( 2B\,(F_C - 1) + 2 B_C \tanh(\delta_C/2) \Big).
\]
\end{tcolorbox}

\emph{Proof.} Run the same argument with the mixed estimators. For $C \in S$,
$e^{b_C} = 1$ and the value error vanishes, so
$A \le 1 + \tilde T$ with $\tilde T = \sum_{C\notin S} \hat p^{\mathrm{mix}}_C (F_C-1)$,
where $\hat p^{\mathrm{mix}}$ uses the mixed normalizers. Since $Z_C \ge \hat Z_C$ on
$S$, the mixed normalizer is larger, hence
$\hat p^{\mathrm{mix}}_C \le \hat p_C$ for $C \notin S$ and $\tilde T \le T$; conclude
with Step~1 and the monotonicity of $x \mapsto x/(1+x)$. The value term runs only over
$C \notin S$. $\hfill\square$

The middle expression is a per-decode-step error certificate computable from
quantities the selection pass already has; the right-hand side is additive per cluster
and needs no normalization across clusters.

\paragraph{Efficient verification}
Per cluster one stores: $\bar k$ (vector) and $\bar v$ (vector, for the approximation
itself); either the coordinate box $(k_{\max}, k_{\min})$ or the radius
$r_C = \max_{i}\|k_i - \bar k\|_2$ (scalar) for
$\delta_C \le \|q\| r_C / \sqrt{d_k}$; the diagonal key variance
$D_C = \mathrm{diag}(\Sigma_C)$ (vector) and the off-diagonal remainder
$\lambda_C = \|\Sigma_C - \mathrm{diag}(\Sigma_C)\|_2$ (scalar) for
\[
\sigma_C^2 \ \le\ \hat\sigma_C^2
= \frac{1}{d_k}\Big( \sum_j q_j^2\, D_{C,j} + \lambda_C\, \|q\|_2^2 \Big);
\]
and $B_C$ (scalar). All query-time work is $O(d_k)$ per cluster.

\paragraph{Remark (variance form of the value term)}
With $Y = e^{x}$ as in Lemma 2,
$\mathbb E|Y - \mathbb E Y| \le \sqrt{\operatorname{Var} Y}$ and
$\mathbb E e^{2x} \le G(2\sigma_C, 2\delta_C)$ give
$\sum_i |\tfrac{1}{|C|} - w_i| \le \sqrt{G(2\sigma_C, 2\delta_C) - 1}$, so one may use
$\min\{\,2\tanh(\delta_C/2),\ \sqrt{G(2\sigma_C, 2\delta_C)-1}\,\}$, which helps when
$\sigma_C$ and $\delta_C$ are both small (before $\tanh$ saturates).
